\chapter{Entropy Depth Lower Bound}

\section{Purpose of EntropyDepth}

EntropyDepth measures the minimum number of admissible refinement steps needed
to remove unresolved entropy. It is the depth invariant associated with the
capacity bounds of the previous chapter.

The point of EntropyDepth is not merely to count runtime. It counts structural
progress inside a specified refinement model. A refinement process may be
algorithmic, logical, combinatorial, spectral, or operational. In every case,
the EntropyDepth question is:

\begin{quote}
How many admissible steps are required before the remaining uncertainty or
defect reaches the terminal threshold?
\end{quote}

This chapter develops the basic lower bound and explains how it is used.

\section{Refinement Trajectories}

Let \(X\) be a state space and let \(\mathcal R\) be a class of admissible
refinement maps. A refinement trajectory is a sequence
\[
x_0,x_1,x_2,\ldots
\]
such that
\[
x_{t+1}=R_t(x_t)
\]
for some \(R_t\in\mathcal R\).

\begin{definition}[Terminal set]
A terminal set is a subset \(T\subseteq X\) whose elements count as resolved,
defect-free, classified, or otherwise finished for the problem under study.
\end{definition}

The terminal condition is part of the model. In some settings, terminal means
that entropy is zero. In others, terminal means that a defect functional
vanishes, that a partition stabilizes, or that a certificate has been
identified.

\section{Definition}

\begin{definition}[EntropyDepth]
Let \(x_0\in X\), let \(\mathcal R\) be an admissible refinement class, and let
\(T\subseteq X\) be a terminal set. The EntropyDepth of \(x_0\) is
\[
\ED_{\mathcal R,T}(x_0)
=
\inf\{t\ge 0:x_t\in T
\text{ for some admissible trajectory }x_0,\ldots,x_t\}.
\]
When \(\mathcal R\) and \(T\) are fixed, we write \(\ED(x_0)\).
\end{definition}

EntropyDepth is infinite if no admissible trajectory reaches the terminal set.

\begin{definition}[Entropy terminal condition]
Given an entropy functional \(H:X\to\mathbb R_{\ge 0}\), the entropy-terminal
set is
\[
T_H=\{x\in X:H(x)=0\}.
\]
\end{definition}

\section{One-Step Entropy Loss}

The depth lower bound depends on the largest entropy loss allowed in one step.

\begin{definition}[Step loss]
For a trajectory \(x_0,x_1,\ldots\), define the step entropy loss
\[
\Delta H_t=H(x_t)-H(x_{t+1}).
\]
\end{definition}

\begin{definition}[Uniform loss bound]
A refinement class \(\mathcal R\) has uniform entropy-loss bound \(c>0\) if
\[
H(x)-H(Rx)\le c
\]
for every \(x\in X\) and every \(R\in\mathcal R\).
\end{definition}

The constant \(c\) is the capacity bound in entropy units. It may come from a
locality estimate, a finite alphabet, a detector bound, a communication
constraint, or another admissibility argument.

\section{Lower Bound}

\begin{theorem}[EntropyDepth Lower Bound]
Let \(H:X\to\mathbb R_{\ge 0}\) be an entropy functional. Suppose every
admissible refinement step removes at most \(c>0\) units of entropy. If
terminal resolution requires \(H(x_t)=0\), then
\[
\ED(x_0)\ge \frac{H(x_0)}{c}.
\]
\end{theorem}

\begin{proof}
Let \(x_0,x_1,\ldots,x_t\) be any admissible trajectory reaching a terminal
state with \(H(x_t)=0\). Since each step removes at most \(c\) entropy,
\[
H(x_i)-H(x_{i+1})\le c
\]
for each \(i<t\). Summing over \(i=0,\ldots,t-1\) gives
\[
H(x_0)-H(x_t)\le tc.
\]
Because \(H(x_t)=0\), this becomes \(H(x_0)\le tc\). Hence
\[
t\ge \frac{H(x_0)}{c}.
\]
Taking the infimum over all admissible terminal trajectories gives the result.
\end{proof}

\section{Linear Entropy Regime}

The most common URF application occurs when initial entropy grows linearly
with system size while the per-step capacity remains bounded.

\begin{theorem}[Linear EntropyDepth Criterion]
Let \(x_n\) be a family of inputs indexed by size \(n\). Suppose there are
constants \(a>0\) and \(c>0\), independent of \(n\), such that
\[
H(x_n)\ge an
\]
and every admissible refinement step removes at most \(c\) entropy. Then
\[
\ED(x_n)=\Omega(n).
\]
\end{theorem}

\begin{proof}
By the EntropyDepth lower bound,
\[
\ED(x_n)\ge \frac{H(x_n)}{c}\ge \frac{a}{c}n.
\]
Thus EntropyDepth grows at least linearly.
\end{proof}

\section{Superlinear and Logarithmic Regimes}

The same calculation also records other growth regimes.

\begin{theorem}[General Growth Transfer]
Let \(g(n)\) be a nonnegative function. If
\[
H(x_n)\ge g(n)
\]
and every admissible step removes at most \(c>0\) entropy, then
\[
\ED(x_n)\ge \frac{g(n)}{c}.
\]
\end{theorem}

\begin{proof}
Apply the EntropyDepth lower bound with \(H(x_n)\ge g(n)\).
\end{proof}

Thus logarithmic entropy gives a logarithmic lower bound, linear entropy gives
a linear lower bound, and superlinear entropy gives a superlinear lower bound,
provided the per-step capacity remains bounded.

\section{Worked Example: Finite Search Space}

Let \(\Omega_n\) be a search space with
\[
|\Omega_n|=2^n.
\]
Using base-two entropy,
\[
H_0=\log_2|\Omega_n|=n.
\]
If each admissible refinement step removes at most \(c\) bits, then every
terminal refinement trajectory has length at least
\[
\frac{n}{c}.
\]

This example is only a refinement-model lower bound. It does not rule out an
algorithm that uses a different admissible model or accesses a certificate not
represented by the refinement process.

\section{Worked Example: Partition Refinement}

Let \(P_t\) be a partition of a configuration space \(\Omega\). Suppose the
true configuration lies in a block \(B_t\in P_t\). Define
\[
H(P_t)=\log |B_t|.
\]
If a refinement step can split \(B_t\) into at most \(B\) effective subblocks,
then the entropy loss per step is at most \(\log B\). Therefore, if
\[
|B_0|=N,
\]
then
\[
\ED(P_0)\ge \frac{\log N}{\log B}.
\]

This is the basic depth estimate for bounded-branching refinement.

\section{Connection to SAT-Style Search}

In a SAT-style search problem with \(n\) Boolean variables, the naive
configuration space has size \(2^n\). The entropy of the full assignment space
is \(n\) bits. A local refinement method that removes only \(O(1)\) bits per
admissible step therefore requires \(\Omega(n)\) refinement depth to isolate a
single assignment.

This statement is conditional on the refinement model. It does not prove
\(P\ne NP\). A complexity-theoretic lower bound would require an independent
argument showing that the relevant class of algorithms normalizes to the
bounded-capacity refinement model.

\section{Connection to Weisfeiler--Lehman Refinement}

In graph refinement, a partition of vertices is updated by local neighborhood
types. If each round separates only bounded local information, EntropyDepth
measures how many rounds are required before the partition stabilizes or
separates the target structure.

For a graph on \(n\) vertices, a refinement process may stabilize in at most
\(n\) strict partition refinements, but the lower-bound question asks whether
some family requires many such steps. EntropyDepth supplies the bookkeeping
language for that question.

\section{Audit Checklist}

A valid EntropyDepth argument should provide:

\begin{center}
\begin{tabular}{ll}
\toprule
Field & Required content \\
\midrule
State space & configurations or refinement states \\
Trajectory & admissible update rule \\
Entropy & functional decreasing toward terminal resolution \\
Step bound & proof that one step removes at most \(c\) entropy \\
Terminal set & condition defining completed refinement \\
Initial growth & lower bound on \(H(x_0)\) or \(H(x_n)\) \\
Conclusion & resulting lower bound on \(\ED\) \\
\bottomrule
\end{tabular}
\end{center}

If the step bound or terminal condition is missing, the EntropyDepth estimate
is not yet established.

\section{Boundary}

This chapter proves the elementary EntropyDepth lower bound and its growth
consequences inside a specified refinement model. It does not prove that all
algorithms are bounded-capacity refinements, does not prove SAT lower bounds,
does not prove graph-isomorphism lower bounds, and does not prove unrestricted
Chronos-RR, unrestricted H4.1/FGL, P vs NP, or any Clay problem.
